<Anticipations of Perception.>

The principle, which anticipates all perceptions, as such,
runs thus: In all appearances the sensation, and the real,
which corresponds to it in the object (realitas
phaenomenon), has an intensive magnitude, i.e., a degree.

<Its principle is: In all appearances the real, which is an object
of the sensation, has intensive magnitude, i.e., a degree.>

<Proof
Perception is empirical consciousness, i.e., one in which there is at the
same time sensation. Appearances, as objects of perception, are not
pure (merely formal) intuitions, like space and time (for these cannot be
perceived in themselves). They therefore also contain in addition to the
intuition the materials for some object in general (through which
something existing in space or time is represented), i.e., the real of the
sensation, as merely subjective representation, by which one can only
be conscious that the subject is affected, and which one relates to an ob-
ject in general. Now from the empirical consciousness to the pure con-
sciousness a gradual alteration is possible, where the real in the former
entirely disappears, and a merely formal (a priori) consciousness of the
manifold in space and time remains; thus there is also possible a syn-
thesis of the generation of the magnitude of a sensation from its begin-
ning, the pure intuition 0, to any arbitrary magnitude.

Now I call that magnitude which can only be apprehended as a unity,
and in which multiplicity can only be represented through approxima-
tion to negation = 0, intensive magnitude. Thus every reality in the
appearance has intensive magnitude, i.e., a degree.>


